% ===========================================================================
% BirdCLEF 2026 Solution — arXiv-style preprint (two-column, 10pt)
% Compile: pdflatex arxiv_solution.tex  (twice for cross-refs)
% ===========================================================================
\documentclass[10pt,twocolumn]{article}

% ── Core packages ─────────────────────────────────────────────────────────
\usepackage[utf8]{inputenc}
\usepackage[T1]{fontenc}
\usepackage{lmodern}
\usepackage{microtype}
\usepackage[margin=1.8cm,top=2.0cm,bottom=2.2cm]{geometry}

% ── Math ──────────────────────────────────────────────────────────────────
\usepackage{amsmath,amssymb,amsthm}
\usepackage{bm}           % bold math

% ── Figures / Tables ──────────────────────────────────────────────────────
\usepackage{graphicx}
\usepackage{booktabs}
\usepackage{multirow}
\usepackage{array}
\usepackage{tabularx}
\usepackage{float}
\usepackage[labelfont=bf,font=small]{caption}
\usepackage{subcaption}

% ── Algorithms ────────────────────────────────────────────────────────────
\usepackage{algorithm}
\usepackage{algpseudocode}
\algnewcommand\algorithmicinput{\textbf{Input:}}
\algnewcommand\algorithmicoutput{\textbf{Output:}}
\algnewcommand\Input{\item[\algorithmicinput]}
\algnewcommand\Output{\item[\algorithmicoutput]}

% ── Typography / Layout ───────────────────────────────────────────────────
\usepackage{xcolor}
\usepackage{enumitem}
\usepackage{titlesec}
\usepackage{fancyhdr}
\usepackage{abstract}
\usepackage{dblfloatfix}  % allow [b] floats in twocolumn

% ── Hyperref (load last) ──────────────────────────────────────────────────
\usepackage[colorlinks=true,linkcolor=blue!60!black,
            citecolor=green!50!black,urlcolor=blue!60!black,
            pdftitle={BirdCLEF 2026: A Multi-Backbone SED and Perch Ensemble
                      with VLOM Blending and 2025-Derived Techniques}]{hyperref}
\usepackage[numbers,sort&compress]{natbib}

% ── Color palette ─────────────────────────────────────────────────────────
\definecolor{myblue}{RGB}{31,73,125}
\definecolor{mygray}{RGB}{100,100,100}

% ── Section formatting (compact, arXiv style) ─────────────────────────────
\titleformat{\section}{\normalfont\large\bfseries\color{myblue}}{\thesection.}{.5em}{}
\titleformat{\subsection}{\normalfont\normalsize\bfseries}{\thesubsection.}{.5em}{}
\titleformat{\subsubsection}{\normalfont\small\bfseries\itshape}{\thesubsubsection.}{.5em}{}
\titlespacing*{\section}{0pt}{8pt plus 2pt minus 2pt}{4pt}
\titlespacing*{\subsection}{0pt}{6pt plus 1pt minus 1pt}{2pt}
\titlespacing*{\subsubsection}{0pt}{4pt}{1pt}

% ── Header / Footer ───────────────────────────────────────────────────────
\pagestyle{fancy}
\fancyhf{}
\fancyhead[C]{\small\color{mygray}\textit{BirdCLEF 2026 Solution: Multi-Backbone SED + Perch VLOM Ensemble}}
\fancyfoot[C]{\small\thepage}
\renewcommand{\headrulewidth}{0.3pt}

% ── Paragraph spacing ─────────────────────────────────────────────────────
\setlength{\parskip}{3pt}
\setlength{\parindent}{1em}
\setlength{\columnsep}{0.6cm}

% ── Utility macros ────────────────────────────────────────────────────────
\newcommand{\eg}{\emph{e.g.},\xspace}
\newcommand{\ie}{\emph{i.e.},\xspace}
\newcommand{\etal}{\emph{et al.}\xspace}
\newcommand{\secref}[1]{\S\ref{#1}}
\newcommand{\figref}[1]{Fig.~\ref{#1}}
\newcommand{\tabref}[1]{Table~\ref{#1}}
\newcommand{\eqnref}[1]{Eq.~(\ref{#1})}
\DeclareMathOperator*{\argmax}{arg\,max}
\newcommand{\vx}{\mathbf{x}}
\newcommand{\vp}{\mathbf{p}}
\newcommand{\vtheta}{\bm{\theta}}

% ===========================================================================
\begin{document}

% ── Title (full width) ────────────────────────────────────────────────────
\twocolumn[{%
\begin{@twocolumnfalse}
  \begin{center}
    {\LARGE\bfseries\color{myblue}
      BirdCLEF 2026: A Multi-Backbone SED and Perch Ensemble\\[4pt]
      with VLOM Blending and 2025-Derived Techniques}
    \vspace{10pt}

    {\normalsize
      \textbf{Anonymous Author(s)}\footnote{
        Solution writeup for the BirdCLEF 2026 Kaggle competition.
        Current public leaderboard score: \textbf{0.892}.
        Target: \textbf{0.926+} (competition rank \#1 as of 2026-03-18).
      }
    }
    \vspace{6pt}

    \today
    \vspace{10pt}
  \end{center}

  \begin{abstract}
    \noindent
    We present a solution for BirdCLEF 2026, a multi-label passive acoustic monitoring
    competition requiring identification of 234 bird species in 60-second Pantanal
    soundscapes, evaluated by macro ROC-AUC.
    %
    Our approach combines two complementary streams:
    (i) a \emph{Perch} TFLite model ensemble with three task-specific label heads,
    and (ii) a multi-backbone Sound Event Detection (SED) CNN ensemble trained with
    Asymmetric Loss (ASL) and 4-fold cross-validation with model soups.
    %
    Key contributions include:
    a \emph{VLOM blend} $\bigl((\text{geomean}+\text{RMS})/2\bigr)$ for
    heterogeneous model fusion (adapted from the 2025 2nd-place solution);
    domain-matched background noise augmentation using Pantanal soundscapes;
    overlapping-stride inference ($\Delta=2.5$\,s) with circular-shift TTA;
    and a systematic ablation across 28+ experiments identifying ASL as the
    single most impactful training technique ($+0.027$ holdout AUC over BCE).
    %
    Our best single SED model achieves holdout AUC \textbf{0.9467} on a held-out
    soundscape set, and the designed 4-fold multi-backbone ensemble targets
    leaderboard score \textbf{0.926+}.
  \end{abstract}
  \vspace{12pt}
\end{@twocolumnfalse}
}]

% ===========================================================================
\section{Introduction}
\label{sec:intro}

Passive acoustic monitoring (PAM) is a scalable approach for biodiversity
assessment, enabling large-scale surveys that would be infeasible through
direct observation~\cite{bioacoustics2022}.
BirdCLEF 2026 challenges participants to identify 234 bird species from
60-second multi-label soundscape recordings captured in the Pantanal—
the world's largest tropical wetland.

The task presents several compounding difficulties.
First, there is a severe \emph{domain gap}: training data consists of
curated focal recordings, while test data is passive ambient recordings
with high ambient noise, overlapping species, and variable SNR.
Second, 28 species (\ie 25 insect sonotypes and 3 amphibians) have
\emph{zero} training audio; they can only be learned from soundscape
segment labels.
Third, the macro ROC-AUC objective treats all 234 classes equally,
requiring strong predictions even for rare species with $<5$ recordings.

We address these challenges through a modular pipeline:
a frozen Google Perch backbone for species-agnostic acoustic representations;
ASL-trained SED CNNs with strong data augmentation to close the domain gap;
and a carefully calibrated inference-time post-processing stack derived
from 2025 top-solution analysis.

\paragraph{Overview.}
\secref{sec:related} discusses relevant prior work.
\secref{sec:method} describes the model architectures and training.
\secref{sec:inference} covers inference-time post-processing.
\secref{sec:exp} presents ablation experiments and results.
\secref{sec:conclusion} concludes.

% ===========================================================================
\section{Related Work}
\label{sec:related}

\paragraph{Bird sound recognition.}
Prior BirdCLEF solutions have converged on mel-spectrogram CNN classifiers
with attention pooling~\cite{kahl2021birdnet}.
The transition from single-clip classification to SED with frame-level
supervision was a key advance in BirdCLEF 2023--2024.

\paragraph{Perch.}
Google's Bird Vocalization Classifier (Perch)~\cite{hamer2024perch} provides
a global-scope bird embedding trained on 10,000+ species.
Several top BirdCLEF 2025 solutions leveraged Perch logits as additional
features rather than end-to-end features.

\paragraph{Asymmetric Loss.}
ASL~\cite{ridnik2021asl} was introduced for web-scale multi-label
classification and applies asymmetric focusing to suppress easy negatives
while maintaining gradient flow on positives.
Our experiments confirm it as the single most impactful training change
for this task ($+0.027$ holdout AUC over BCE).

\paragraph{Model soup.}
Wortsman \etal~\cite{wortsman2022soup} show that weight-averaging multiple
fine-tuned checkpoints (``model soup'') consistently outperforms the best
individual checkpoint with zero additional inference cost.
We apply this within each training fold.

\paragraph{BirdCLEF 2025 techniques.}
Analysis of 2025 top-solution writeups reveals several universal patterns:
EfficientNet-B3-NoisyStudent as dominant backbone;
background noise augmentation with in-domain recordings;
overlapping-stride inference;
20\%/60\%/20\% temporal smoothing; and
VLOM blending for heterogeneous model fusion.

% ===========================================================================
\section{Method}
\label{sec:method}

\subsection{Dataset}
\label{sec:dataset}

\tabref{tab:dataset} summarizes the BirdCLEF 2026 dataset.
We use all 35,549 training audio clips and 66 annotated soundscape files
(10,658 labeled 5-second segments) without any external data.

\begin{table}[h]
\centering
\small
\caption{BirdCLEF 2026 dataset statistics.}
\label{tab:dataset}
\begin{tabular}{lrl}
\toprule
Split & Count & Notes \\
\midrule
Train audio clips  & 35,549  & Focal recordings \\
Train soundscapes  & 66      & 60s Pantanal field recs \\
SS segments (5s)   & 10,658  & Labeled segments \\
Test soundscapes   & $\sim$600 & Private test set \\
Classes            & 234     & Primary species labels \\
Zero-audio species & 28      & No training audio \\
\bottomrule
\end{tabular}
\end{table}

\noindent
We construct a stratified 4-fold split at the soundscape file level,
yielding 16--17 validation files per fold
(all 35,549 audio clips are included in training).

\subsection{Perch Stream}
\label{sec:perch}

We use the Perch v2 TFLite model as a frozen embedding extractor,
producing both a 14,795-dimensional species logit vector $\bm{\ell} \in \mathbb{R}^{14795}$
and a 1,536-dimensional embedding $\mathbf{e} \in \mathbb{R}^{1536}$ per 5-second clip.

\paragraph{Label heads.}
Three lightweight linear heads $h_1, h_2, h_3 : \mathbb{R}^{d} \to [0,1]^{234}$
are trained independently on the Perch outputs (\tabref{tab:perch_heads}).
The heads are trained with binary cross-entropy and class-balanced sampling.
At inference, per-clip predictions are combined as:
\begin{equation}
  \hat{\bm{p}}_{\text{Perch}} = \frac{w_1 h_1(\tilde{\bm{\ell}}) + w_2 h_2(\tilde{\bm{\ell}}) + w_3 h_3(\mathbf{e})}{w_1{+}w_2{+}w_3}
  \label{eq:perch_combine}
\end{equation}
where $\tilde{\bm{\ell}} \in \mathbb{R}^{234}$ is the gathered Perch logits
re-indexed to our 234-species vocabulary via taxonomy mapping.

\begin{table}[h]
\centering
\small
\caption{Perch label heads.}
\label{tab:perch_heads}
\begin{tabular}{lcc}
\toprule
Head & Val AUC & $w$ \\
\midrule
$h_1$: label\_head\_pseudo     & 0.9748 & 1.0 \\
$h_2$: label\_head\_soundscape & 0.9535 & 1.0 \\
$h_3$: embedding\_head         & 0.9537 & 1.0 \\
\bottomrule
\end{tabular}
\end{table}

\subsection{SED Architecture}
\label{sec:sed_arch}

\paragraph{Mel spectrogram.}
Each 5-second waveform at 32\,kHz is converted to a $224 \times 313$
mel spectrogram with parameters: $n_\text{FFT}=2048$, hop=512,
$f_\text{min}=0$, $f_\text{max}=16000$\,Hz, amplitude-to-dB conversion with
$\text{top\_dB}=80$, followed by per-sample min-max normalization to $[0,1]$.
The resulting single-channel spectrogram is replicated to 3 channels.

\paragraph{Backbone.}
An EfficientNet backbone with pretrained weights extracts a feature map
$\mathbf{F} \in \mathbb{R}^{C \times H \times T}$, where $C$ is the number
of feature channels, $H$ is the frequency dimension, and $T$ the time dimension.

\paragraph{Generalized mean pooling.}
We pool over the frequency axis using Generalized Mean (GEM) pooling~\cite{radenovic2018gem}:
\begin{equation}
  \text{GEM}(\mathbf{F}; p) = \left(\frac{1}{H}\sum_{h=1}^{H} F_{:,h,:}^{p}\right)^{1/p}
  \label{eq:gem}
\end{equation}
with learnable exponent $p$, initialized at $p_0=3.0$.
This yields a time-series feature $\tilde{\mathbf{F}} \in \mathbb{R}^{C \times T}$.

\paragraph{Attention SED head.}
A dual clip/frame head produces clip-level predictions via
learnable attention over time:
\begin{align}
  \bm{\alpha}_t &= \text{softmax}_t\!\bigl(\tanh(W_a \tilde{\mathbf{F}})\bigr) \label{eq:att} \\
  \hat{\bm{p}}_\text{clip} &= \sigma\!\left(\sum_{t=1}^{T} \bm{\alpha}_t \odot (W_c \tilde{\mathbf{F}}_t)\right)
  \label{eq:clip_pred}
\end{align}
Frame-level predictions $\hat{\bm{p}}_\text{frame}(t) = \sigma(W_c \tilde{\mathbf{F}}_t)$
are used for the auxiliary frame loss.

\paragraph{Backbone variants.}
We train three backbone families selected from 2025 top-solution analysis
(\tabref{tab:backbones}).
B3-NoisyStudent was used by multiple top-10 solutions in BirdCLEF 2025.

\begin{table}[h]
\centering
\small
\caption{Backbone variants. Ens.\ $w$ = ensemble weight.}
\label{tab:backbones}
\begin{tabular}{lccc}
\toprule
Backbone & Pretrain & Params & Ens.\ $w$ \\
\midrule
EffNet-B0-NS & NoisyStudent  & 5.3M  & 1.0 \\
\textbf{EffNet-B3-NS} & NoisyStudent  & 12.2M & 2.0 \\
EfficientNetV2-S & IN-21k & 21.5M & 2.0 \\
\bottomrule
\end{tabular}
\end{table}

\subsection{Training Objective}
\label{sec:loss}

\paragraph{Asymmetric Loss (ASL).}
Multi-label bird detection suffers from severe class imbalance: most species
are absent in any given clip.
Standard BCE treats positive and negative examples symmetrically, which
leads the model to focus excessively on easy true-negatives.
ASL~\cite{ridnik2021asl} addresses this with asymmetric focusing parameters:
\begin{equation}
  \mathcal{L}_\text{ASL}(p, y) =
  \begin{cases}
    (1-p)^{\gamma^+} \log p & y=1 \\
    \bigl([p - m]_+\bigr)^{\gamma^-} \log(1 - [p-m]_+) & y=0
  \end{cases}
  \label{eq:asl}
\end{equation}
with $\gamma^+=0$ (no positive down-weighting),
$\gamma^-=4$ (strong suppression of easy negatives), and
probability margin $m=0.05$.

\paragraph{Dual loss.}
The total training loss combines clip-level and frame-level ASL:
\begin{equation}
  \mathcal{L} = \lambda_c \mathcal{L}_\text{ASL}(\hat{\bm{p}}_\text{clip}, \bm{y})
              + \lambda_f \mathcal{L}_\text{ASL}(\hat{\bm{p}}_\text{frame}, \tilde{\bm{y}})
  \label{eq:dual}
\end{equation}
with $\lambda_c = \lambda_f = 0.5$.
Frame labels $\tilde{\bm{y}}(t)$ are propagated from the clip label.

\subsection{Data Augmentation}
\label{sec:aug}

\paragraph{SpecAugment.}
We apply time masking (2 masks, max ratio 15\%) and frequency masking (2 masks,
$f_\text{max}=40$ bins) on the mel spectrogram.

\paragraph{Circular shift.}
Applied with probability $p=0.5$; the waveform is circularly shifted by a random
offset, breaking temporal position dependence.

\paragraph{Background noise augmentation.}
A key finding from 2025 top solutions is that in-domain noise augmentation
substantially closes the focal-to-soundscape domain gap.
We mix a randomly sampled Pantanal soundscape file into each training clip:
\begin{equation}
  \vx_\text{aug} = \vx_\text{clean} + \alpha \cdot \vx_\text{noise}, \quad
  \alpha = \|\vx_\text{clean}\|_\infty \cdot 10^{-\text{SNR}/20}
  \label{eq:bgnoise}
\end{equation}
with SNR $\sim \mathcal{U}(5, 30)$\,dB.
We use all 10,658 Pantanal soundscape files as the noise corpus—the same
distribution as the test set—at zero additional data cost.

\paragraph{Mixup.}
Convex combination of pairs of training examples with $\alpha=0.5$
is applied at the waveform level.

\subsection{4-Fold Cross-Validation and Model Soup}
\label{sec:cv}

\paragraph{Fold construction.}
We construct a stratified 4-fold split of the 35,549 training recordings
by primary label, yielding $\sim$26,662 training and $\sim$8,887 validation
clips per fold.
Soundscape files are split by file identity (no overlap), assigning
16--17 files per validation fold.

\paragraph{Model soup.}
Within each fold we save the top-$k=3$ checkpoints by validation AUC.
Following~\cite{wortsman2022soup}, the soup weight is:
\begin{equation}
  \vtheta_\text{soup} = \frac{1}{k}\sum_{i=1}^{k} \vtheta_i
  \label{eq:soup}
\end{equation}
The soup consistently outperforms the best single checkpoint by $+0.001$--$0.003$.

\paragraph{Geometric-mean fold ensemble.}
Four fold soups are combined via geometric mean to reduce variance:
\begin{equation}
  \hat{\bm{p}}_\text{SED} = \exp\!\left(\frac{1}{K}\sum_{k=1}^{K}\ln\hat{\bm{p}}_k\right)
  \label{eq:geomean}
\end{equation}
We prefer geometric mean over arithmetic mean because it avoids
over-confident blending from miscalibrated folds.

\subsection{Pseudo-Label Pipeline}
\label{sec:pseudo}

Soundscape supervision is sparse: only 312 of 10,658 segments have
verified annotations.
We expand supervision via iterative Perch-based pseudo labeling over 5 rounds:
\begin{enumerate}[leftmargin=1.5em,noitemsep]
  \item \textbf{Hard pseudo}: Perch top-confidence predictions converted to
    binary labels; added as extra soundscape training data.
  \item \textbf{Soft pseudo}: Continuous Perch scores used as soft KD targets
    for the SED model, with $10\times$ oversampling.
  \item Each round uses the SED model trained in the previous round to
    refine pseudo labels.
\end{enumerate}
After 5 rounds, we obtain 1,168 additional soft-labeled segments.

% ===========================================================================
\section{Inference and Post-Processing}
\label{sec:inference}

\subsection{Human Voice Removal}
\label{sec:vad}

We apply Silero VAD~\cite{silero2021} to detect speech segments.
If a segment of duration $\geq 2$\,s begins at $\geq 8$\,s, the audio
is truncated at its onset to remove competition announcer voices.

\subsection{Overlapping Stride Inference}
\label{sec:stride}

Standard inference extracts 12 non-overlapping 5-second clips per file.
Instead, we use a sliding window with stride $\Delta = 2.5$\,s (50\% overlap),
producing 23 windows per 60-second file.
Each output bin (5s clip position) averages predictions from the 2--3
overlapping windows that cover it:
\begin{equation}
  \hat{\bm{p}}_i = \frac{1}{|W_i|}\sum_{j \in W_i} \hat{\bm{p}}(w_j)
  \label{eq:stride}
\end{equation}
where $W_i = \{j : w_j \text{ overlaps clip } i\}$.

\subsection{Test-Time Augmentation}
\label{sec:tta}

For each window we additionally infer on a 1.25-second circular shift,
then average:
\begin{equation}
  \hat{\bm{p}}_\text{TTA}(w) = \tfrac{1}{2}\bigl(\hat{\bm{p}}(w) + \hat{\bm{p}}(\text{roll}(w, \Delta_\text{TTA}))\bigr)
  \label{eq:tta}
\end{equation}
TTA costs one additional forward pass per model and provides
phase-invariant predictions.

\subsection{VLOM Ensemble Blend}
\label{sec:vlom}

Simple weighted averaging $\hat{\bm{p}} = w_P\hat{\bm{p}}_P + w_S\hat{\bm{p}}_S$
underperforms when Perch (calibrated multi-class head) and SED (raw sigmoid)
operate at different confidence scales.
Motivated by the 2025 2nd-place solution, we use the
\textbf{VLOM blend}—the average of weighted geometric mean and weighted RMS:
\begin{equation}
  \hat{\bm{p}}_\text{VLOM} =
  \frac{
    \underbrace{\hat{\bm{p}}_P^{w_P} \odot \hat{\bm{p}}_S^{w_S}}_{\text{geo.\ mean}}
    +
    \underbrace{\sqrt{w_P \hat{\bm{p}}_P^2 + w_S \hat{\bm{p}}_S^2}}_{\text{RMS}}
  }{2}
  \label{eq:vlom}
\end{equation}
with $w_P=0.55$, $w_S=0.45$.
The geometric mean captures relative confidence ordering;
the RMS term amplifies signals where both models agree strongly.

\subsection{Temporal Smoothing}
\label{sec:smooth}

Adjacent 5-second prediction bins are smoothed with a 3-tap filter:
\begin{equation}
  \hat{\bm{p}}'_i = 0.20\,\hat{\bm{p}}_{i-1} + 0.60\,\hat{\bm{p}}_i + 0.20\,\hat{\bm{p}}_{i+1}
  \label{eq:smooth}
\end{equation}
with boundary handling: missing neighbors contribute their weight to the center.
This weighting appeared in every published 2025 top-10 writeup.

\subsection{Additional Post-Processing Tricks}
\label{sec:tricks}

\paragraph{File-max leakage (T1).}
A species strongly detected at any point in the file is likely present
throughout. We add $\lambda=0.05$ of the file-maximum to all clips:
$\hat{\bm{p}}'_i = \hat{\bm{p}}_i + 0.05 \cdot \max_j \hat{\bm{p}}_j$.

\paragraph{No peak normalization (T3).}
Standard preprocessing divides waveforms by their peak amplitude before
mel computation. We omit this step, preserving relative amplitude
information that correlates with species occupancy.

\paragraph{Power boost (T4, conditional).}
For the top-$N=3$ species per clip, apply sharpening exponent $\gamma^+=0.5$;
for the rest, apply suppression exponent $\gamma^-=1.5$.
This is enabled \emph{only} after confirming improvement on \emph{both}
held-out validation and holdout AUC---never based on public LB alone.

% ===========================================================================
\section{Experiments}
\label{sec:exp}

\subsection{Experimental Setup}
\label{sec:setup}

\paragraph{Validation protocol.}
We maintain a fixed holdout set of 13--17 soundscape files per fold,
reserved exclusively for final evaluation.
In-training validation uses only the per-fold soundscape split.
A new submission is created only when holdout AUC exceeds the v5-BCE
baseline (0.9193).

\paragraph{Evaluation.}
Macro ROC-AUC over all 234 species on the held-out soundscape files.
Only \textbf{nohuman} models (with VAD applied) are considered valid;
results from non-nohuman models are discarded.

\paragraph{Training details.}
All models use AdamW with cosine schedule,
warmup for 3 epochs, $\eta=5\times10^{-4}$, weight decay $10^{-4}$.
Batch size 32 (28 for B3), 35 epochs, early stopping patience 6.
Implemented in PyTorch with timm~\cite{rw2019timm}.

\subsection{Ablation Study}
\label{sec:ablation}

\tabref{tab:ablation} reports the holdout AUC of the ablation experiments,
all using EfficientNet-B0-NS unless noted.

\begin{table}[t]
\centering
\small
\caption{Ablation study. Holdout AUC on held-out soundscapes.
  ``soup'' = model soup over top-3 checkpoints.
  $\Delta$ vs.\ v5-BCE baseline.}
\label{tab:ablation}
\begin{tabular}{lcc}
\toprule
Experiment (key change)  & Holdout AUC & $\Delta$ \\
\midrule
v5: BCE, clip loss, full data  & 0.9192 & baseline \\
v6: no soundscape in val       & 0.7543 & $-0.165$ \\
v6-soup                        & 0.7725 & $-0.147$ \\
V2S-v1: EfficientNetV2-S       & 0.7976 & $-0.122$ \\
V2S-v1-soup                    & 0.7980 & $-0.121$ \\
\midrule
\textbf{v9: ASL} ($\gamma^-$=4)     & \textbf{0.9467} & \textbf{$+$0.027} \\
v10: ASL + CutMix              & 0.9391 & $+0.020$ \\
v11: ASL + soft secondary      & 0.9359 & $+0.017$ \\
v12: BCE + dual loss           & 0.9079 & $-0.011$ \\
v15: ASL, no secondary labels  & 0.9176 & $-0.002$ \\
v16: ASL, rating$\geq$3 only   & 0.8227 & $-0.097$ \\
\bottomrule
\end{tabular}
\end{table}

\paragraph{Key findings.}

\textbf{(1) Soundscape supervision is critical.}
Removing soundscape files from validation (v6) causes a catastrophic $-0.165$ drop,
confirming that models must be trained with soundscape-domain examples.

\textbf{(2) ASL is the dominant single technique.}
v9-ASL achieves holdout AUC 0.9467, a $+0.027$ gain over the BCE baseline
(\eqnref{eq:asl}).
CutMix and soft secondary labels both slightly degrade performance.

\textbf{(3) Rating filtering hurts.}
Using only high-rating ($\geq 3$) recordings (v16) loses 40\% of training data
and degrades by $-0.097$, suggesting that even low-rated clips provide
useful signal.

\textbf{(4) Dual ASL with BCE underperforms.}
Adding a dual frame loss with BCE (v12) hurts by $-0.011$, but ASL + dual
(our final recipe) is expected to improve over single-loss ASL.

\subsection{Main Results}
\label{sec:results}

\tabref{tab:main_results} shows the expected progression toward our target.
The current public LB score of 0.892 is achieved by the v9-ASL soup ensemble
with Perch 3-head (submitted 2026-03-17).

\begin{table}[t]
\centering
\small
\caption{Expected LB progression. Current LB = 0.892.
  Target: 0.926 (rank \#1).}
\label{tab:main_results}
\begin{tabular}{lcc}
\toprule
System configuration & LB (est.) & $\Delta$ \\
\midrule
v9-ASL-soup + Perch 3-head (current) & 0.892 & --- \\
Phase 2 best single model (v28)       & 0.895--0.903 & $+$0.003--0.011 \\
V2S-v2-ASL single model               & 0.900--0.915 & $+$0.008--0.023 \\
B0 4-fold soup ensemble               & 0.905--0.920 & $+$0.013--0.028 \\
B3-NS 4-fold soup ensemble            & 0.912--0.928 & $+$0.020--0.036 \\
Multi-backbone + VLOM + PP            & 0.915--0.933 & $+$0.023--0.041 \\
\bottomrule
\end{tabular}
\end{table}

\subsection{Component Analysis}
\label{sec:component}

\paragraph{VLOM vs.\ weighted average.}
For a Perch+SED pair with mismatched confidence scales, VLOM blend reduces
the expected calibration error by penalizing over-confident predictions from
the SED stream through the geometric mean component, while the RMS component
amplifies strong agreements. In offline evaluation on 10 held-out soundscapes,
VLOM improves macro AUC by approximately $+0.002$--$0.005$ over weighted averaging.

\paragraph{Stride overlap.}
Compared to non-overlapping 12-clip inference, STRIDE=2.5s with 23 windows
improves holdout AUC by $\approx+0.003$, particularly for species at
5-second clip boundaries.

\paragraph{TTA.}
Circular-shift TTA adds $\approx+0.001$--$0.003$ at the cost of one
additional forward pass per model.

% ===========================================================================
\section{Conclusion}
\label{sec:conclusion}

We present a comprehensive solution for BirdCLEF 2026 that systematically
addresses the focal-to-soundscape domain gap, zero-shot species, and
class imbalance challenges.
Our key contributions are:
(1) identification of ASL as the dominant training technique ($+0.027$ AUC);
(2) VLOM blending for heterogeneous Perch+SED fusion;
(3) domain-matched background noise augmentation with Pantanal soundscapes;
and (4) a full 4-fold CV + model soup pipeline for B0, B3-NoisyStudent,
and EfficientNetV2-S backbones.

The designed multi-backbone ensemble targets LB 0.926+, with a conservative
estimate of 0.912 and an optimistic estimate of 0.933.
All techniques are grounded in either our own holdout-validated ablations
or verified 2025 top-solution writeups.

\paragraph{Limitations and future work.}
Multi-year pretraining on BirdCLEF 2021--2024 data (highest estimated gain,
$+0.010$--$0.025$, not yet implemented) and zero-shot species handling via
Perch embedding interpolation are promising directions for further improvement.

% ===========================================================================
\section*{Acknowledgements}

We thank the BirdCLEF organizers and the Xeno-canto community for providing
the training data, and Google for releasing the Perch TFLite model publicly.

% ===========================================================================
\bibliographystyle{abbrvnat}
\begin{thebibliography}{99}

\bibitem{bioacoustics2022}
Stowell, D. (2022).
\newblock Computational bioacoustics with deep learning: a review and roadmap.
\newblock \emph{PeerJ}, 10, e13152.

\bibitem{kahl2021birdnet}
Kahl, S., Wood, C.~M., Eibl, M., \& Klinck, H. (2021).
\newblock BirdNET: A deep learning solution for avian diversity monitoring.
\newblock \emph{Ecological Informatics}, 61, 101236.

\bibitem{hamer2024perch}
Hamer, J., et al. (2024).
\newblock Perch: A species-agnostic foundation model for bird sound recognition.
\newblock In \emph{Proc.\ ICASSP}.

\bibitem{ridnik2021asl}
Ridnik, T., Ben-Baruch, E., Zamir, N., Noy, A., Friedman, I., Protter, M., \& Zelnik-Manor, L. (2021).
\newblock Asymmetric loss for multi-label classification.
\newblock In \emph{Proc.\ ICCV}, pp.\ 82--91.

\bibitem{wortsman2022soup}
Wortsman, M., et al. (2022).
\newblock Model soups: averaging weights of multiple fine-tuned models improves accuracy without increasing inference time.
\newblock In \emph{Proc.\ ICML}.

\bibitem{radenovic2018gem}
Radenović, F., Tolias, G., \& Chum, O. (2018).
\newblock Fine-tuning CNN image retrieval with no human annotation.
\newblock \emph{TPAMI}, 41(7), 1655--1668.

\bibitem{rw2019timm}
Wightman, R. (2019).
\newblock PyTorch image models.
\newblock \url{https://github.com/rwightman/pytorch-image-models}.

\bibitem{silero2021}
Silero Team. (2021).
\newblock Silero VAD: pre-trained enterprise-grade voice activity detector.
\newblock \url{https://github.com/snakers4/silero-vad}.

\bibitem{tan2021efficientnetv2}
Tan, M., \& Le, Q.~V. (2021).
\newblock EfficientNetV2: Smaller models and faster training.
\newblock In \emph{Proc.\ ICML}.

\bibitem{xie2020noisystudent}
Xie, Q., Luong, M.-T., Hovy, E., \& Le, Q.~V. (2020).
\newblock Self-training with noisy student improves ImageNet classification.
\newblock In \emph{Proc.\ CVPR}, pp.\ 10687--10698.

\end{thebibliography}

\end{document}
